\scanpage{55}
\mspage{11}
% archive sheet number: 54
3) \deleted{Gommage} \inserted{Bouchage de trous, ou ``gommage''} des cercles de recollement.

Cette opération est relative à la donnée d'une 1-sous-variété fermée de $\partial X$, i.e. d'une partie $A'_{\ell}$ de $\pi_0(\partial X)\simeq A_{\ell}$ (ens des arêtes libres du graphe). Elle consiste à boucher les trous \unsure{correspondants|correspondantes}, au moyen des cônes sur les composantes connexes du bord de $\Sigma'_{\ell}$.

Au niveau du graphe, cela consiste à gommer les arêtes libres correspondantes, avec bien entendu aussi les repères de marquage qui se trouvent au dessus --- \deleted{mais} les pondérations des sommets restent inchangées.

\deleted{\unsure{4|+}}\quad \uline{Gommage simultané mixte}, \deleted{\unsure{au}} correspondant à une partie \deleted{\unsure{quelconque des}} \inserted{$R'\amalg\Sigma'$} \deleted{$A'\amalg R'$} de $R\amalg\Sigma$ qui soit à la fois ouverte et fermée, \deleted{\unsure{i.e.}} \inserted{et telle} que \deleted{$\unsure{\Sigma\amalg A\amalg R|A\amalg R}$} \inserted{$R'\supset R\cap\Sigma'$}, \deleted{un ens. de composantes connexes de $R\amalg\Sigma$}, ou ce qui revient au $=$, une partie \inserted{$R'\amalg A'$} de $R\amalg A$, telle que $R'\supset R|A'$.

Au niveau du graphe pondéré à repères, cette opération consiste à enlever \inserted{de $\vec R$} \deleted{\unsure{les}} la partie corr. $\vec R'$, à contracter en un point les arêtes de $\unsure{A'|\Gamma'}\deleted{\unsure{{}_{\mathrm{pond}}}}$ (aussi bien libres que liées), en faisant suivre la pondération.
